\section{Woran die Idee sich prüfen lässt}
\label{sec:tests}

Die Aufgabe eines gemeinsamen Trägers ist erst dann erfüllt, wenn dieselben Zustände, Symmetrien und Kopplungen in den aufeinanderfolgenden Rechnungen wieder auftauchen. Die Beispiele machen diese Anforderung an drei Stellen konkret. Sie zeigen auch, wie sich ein negativer Befund zur jeweils benutzten Annahme zurückverfolgen lässt.

\subsection{Was die Feldauslese erhält}

Beim HC-Feldvergleich sind die Eingaben eine Feldkonfiguration und die deklarierte Involution \(J_{\mathrm B}\). Abbildung \ref{fig:hc-feld} zeigt den Test am lichtartigen Eingang. Das berechnete \(s'\) sagt, welche Feldklasse nach der Projektion erreicht wird. Ein Modell, das aus demselben Projektor wieder ein lichtartiges Feld gewinnen soll, muss angeben, durch welche anschließende Operation die beiden Bedingungen \(s'=0\) und \(\mathbf z'\ne0\) erfüllt werden. Die Feldinvariante macht diese Forderung unabhängig von der Wahl eines Lorentz-Beobachters prüfbar. Die Identifikation eines quantisierten Lichtzustands setzt weitere Struktur voraus \cite{feldkern,feldsektoren}.

Auch die Dynamik der Auslese hat einen eigenen Test. Für ein vorgegebenes Paar \((T(t),\allowbreak G(t))\) genügt ein \(u_{\ker}\in\ker T\) mit \((\dot T+TG)u_{\ker}\ne0\), um eine autonome lineare Entwicklung der ausgelesenen Variablen auszuschließen. Die Zustände \(u\) und \(u+u_{\ker}\) ergeben dieselbe Auslese und verschiedene Ableitungen. Die Gleichheit der momentanen Beobachtungsdaten entscheidet so noch nicht, ob zwei Zustände dieselbe Antwort auf eine Störung haben. Diese Prüfung lässt sich an einer konkreten Realisierung durchführen, sobald ihre Auslese und ihr Generator feststehen \cite{physics}.

\subsection{Was bei den Massen getrennt festgelegt werden muss}

Im Gram-Vergleich sind die Beträge der paarweisen Überlappungen und die Phase getrennte Eingaben. Man kann die Phase ändern und alle Beträge festhalten. Die Normrelation bleibt erhalten, während die Determinante und das Spektrum sich ändern. Abbildung \ref{fig:koide-phase} zeigt den Rangverlust an einer konkreten Stelle dieser Familie. Der Versuch, die drei Leptonenmassen aus der Überlappung allein zu bestimmen, lässt daher einen ausrechenbaren Teil des Spektrums frei.

Welche Prüfung folgt daraus? Eine vorgeschlagene Zustandspräparation muss zuerst unabhängig die zugehörige Gram-Matrix liefern. Falls sie Rang zwei hat, liefert diese Wurzelregel eine verschwindende Masse. Den Randfall kann man angeben, bevor man eine Skala wählt. Andererseits liefert ein positiv definites Beispiel zunächst drei positive Wurzeln. Für einen Vergleich mit physikalischen Massen braucht man danach die Identifikation des nativen Operators, die Auswahl der Phase und eine unabhängig bestimmte Normierung. Ein erfolgreicher Test der Normrelation beantwortet in diesem Aufbau genau eine der drei Fragen \cite{physics,companion}.

\subsection{Was ein OP5-Vergleich messen müsste}

Bei OP5 ist die Eingabe das präparierte Profil aus \eqref{eq:op5-profile}. Abbildung \ref{fig:op5-modell} zeigt, wie daran zwei Antwortwege anschließen. Man kann die Erhaltung des mechanischen Tensors prüfen und für den gerichteten Weg den Gezeitenkontrast auf dem Plateau berechnen. Der gewöhnliche Außenbeitrag \eqref{eq:op5-compact} und der Achsenkontrast \eqref{eq:op5-contrast} sind damit am selben Profil getrennt zugänglich. Das Nahachsenverhalten in Abbildung \ref{fig:op5-profil} zeigt außerdem, warum ein auf der Achse verschwindender Profilwert den Gezeitentest offenlässt. Die Auslese benutzt Ableitungen des Profils.

Eine spätere Messung braucht für beide Antwortwege einen Bezug zur realen Präparation. Die Beiträge des Materials, der realen Halterung und der Ansteuerungsenergie bestimmen den gewöhnlichen Hintergrund. Die Empfindlichkeit der Anschlüsse und die effektive Kopplungsdifferenz bestimmen, ob sich der gerichtete Kontrast in einer Übertragungsdifferenz wiederfindet. Die integrierte Spur in \eqref{eq:op5-second} lokalisiert ein konkretes Problem der linearen Kompensation in der nächsten Materialordnung. Für diese Ordnung ist ein vollständiger Tensor am räumlichen Ort der Messung zu bilanzieren \cite{op5-elastic-limit,op5-budget}. Der angegebene Grenzfall liefert damit eine überprüfbare Anforderung an eine Apparatur.

\subsection{Was an einer einzigen Zahl hängt}
\label{sec:one-number}

Die Zahl stammt aus der Zählung der Zustände eines Horizonts in der
Schleifenquantengravitation. Eine Horizont-Zählregel ist eine Gleichung
\[
 \sum_j g(j)\,e^{-2\pi\gamma\sqrt{j(j+1)}}=1,\qquad j=\tfrac12,1,\tfrac32,\dots,
\]
deren Lösung \(\gamma\) den Barbero-Immirzi-Parameter festlegt; \(g(j)\) ist das
Gewicht des Spins \(j\). Der Materie-Frame M zählt jeden Spin mit dem Gewicht
\(g(j)=2\), nach der Regel von Domagala und Lewandowski. Das ergibt \(\gamma_M=0{,}237533\). Der
Tachyon-Frame T zählt nur \(j=\tfrac12\), ebenfalls mit dem Gewicht 2, nach der
Regel von Ashtekar, Baez, Corichi und Krasnov. Das ergibt
\(\gamma_T=\ln2/(\pi\sqrt3)=0{,}127384\) (Runden v4.10.0-r52 und v4.10.0-r53
\cite{layers}). Der Quotient ist der Frame-Faktor
\(\gamma_L=\gamma_M/\gamma_T=1{,}8647\). Schreibt man \(\gamma_L=\cosh2\eta\)
und \(\beta=\tanh\eta\), gehört dazu die Geschwindigkeit \(\beta=0{,}5494\)
(Rapidität \(\eta=0{,}6175\)). OP1, das erste offene Problem des Companion,
ist die Frage, warum \(\beta\) diesen Wert hat; gleichwertig: warum M und T
nach verschiedenen Regeln zählen. In der Beobachterlesart ist M ein
Beobachter, der sich relativ zum Ruhesystem eines Feldes mit parallelem
elektrischem und magnetischem Teil bewegt. OP1 fragt dann, warum er sich mit
\(\beta=0{,}5494\) bewegt; \(\gamma_L\) ist in dieser Lesart das Verhältnis
der Energiedichten des Feldes, gemessen von M und im Ruhesystem (Runde
v4.10.0-r29 \cite{companion}).

Seit dem 26. November 2026 ist die Zählregel als \emph{Postulat} angenommen (Runde v4.10.0-r67 \cite{countingpostulate}), im selben Sinn
wie die Gleichberechtigung der Frames: benannt, nicht hergeleitet. Es hat drei Teile. Im Materie-Frame zählt jeder Spin mit dem Gewicht 2
(die beiden äußersten Zustände \(m=\pm j\)); die Fläche eines Durchstoßpunkts ist proportional zu \(\sqrt{j(j+1)}\); und der
Tachyon-Frame zählt nur \(j=\tfrac12\), wobei ein solcher Punkt ein Bit trägt. Mit \(k=2j\) folgt daraus eine einzige Gleichung für den
Frame-Faktor,
\[
 \sum_{k\ge1}2^{\,1-\gamma_L\sqrt{k(k+2)/3}}=1,\qquad \gamma_L=1{,}8646997646,\qquad \beta=0{,}5494056410 .
\]
Damit gibt es eine Rechnung, die zeigt, woher die Zahl kommt; es fehlt der Grund, warum die Zählung diese Form hat. Ersetzt man die Fläche
durch \(j+\tfrac12\), gibt dieselbe Regel exakt \(\beta=1/\sqrt3\). Der Wert ist nicht an Daten angepasst, die Regel wurde aber in
Kenntnis der Datenfenster gewählt. OP1 bleibt offen und lautet jetzt: Leite das Zählpostulat her.

Dieselbe Runde und Runde v4.10.0-r66 halten fest, was das Bild zweier Auslesen dazu beiträgt. Liest man die Hopf-Sphäre als Himmelskugel
von Lichtrichtungen (eine gesetzte Lesart), dann ist \(\beta\) der Überlapp des Hopf-Zustands mit dem Gegenpol des Photonpols, und der
Boost, der beide gegenüberstellt, hat die Geschwindigkeit \(\beta\); Winkel-Lesart und Geschwindigkeits-Lesart sind dann dieselbe
Aussage. Das gilt für jeden Wert von \(\beta\) und leitet keinen her. Die naheliegende Schließungsbedingung dieses Bildes, dass die
Übergangswahrscheinlichkeiten zwischen den zwei Auslesen über alle Spins zusammen eins ergeben, lautet \(\sum_j2\beta^{4j}=1\) und
erzwingt \(\beta=1/\sqrt3\), nicht den gemessenen Wert (Test H165, widerlegt, überwiegend nicht blind). Eine Sichtung zählt 37 frühere
Versuche, die Zahl herzuleiten; keiner gelingt.

Das Zeichen \(\gamma\) hat in diesem Abschnitt drei Bedeutungen. Erstens ist
\(\gamma_L\) ein Frame-Faktor; er hat die Form eines Lorentz-Faktors,
\(\cosh2\eta\), und ist größer als eins. Zweitens sind \(\gamma_M\),
\(\gamma_T\) und das \(\gamma\) des Holst-Terms unten Werte des
Barbero-Immirzi-Parameters. Drittens ist derselbe Wert in geeigneten Einheiten
die Dimension einer selbstähnlichen Menge; \(\gamma_1\) am Ende des Abschnitts
ist eine solche Dimension.

Die bedingten Zahlenwerte des Companion gehen auf diese eine
dimensionslose Zahl zurück. In die
offenen Probleme gelangt sie über fünf getrennte Gleichsetzungen: einen
Spektralversatz, die Restbiegung der Leptonphase, \(|U_{e2}|^2=\beta^2\) in
der Neutrinomischung, \(\Omega_M=\beta^2\) für den Materieanteil und
\(\cot\theta_W=\gamma_L\). Vier davon wurden erst gewählt, als die Daten schon
bekannt waren (Bemerkung v4.10.0-r30.6 \cite{companion}). Jedes Problem behält daneben einen eigenen Rest, der von der
Zahl nicht abhängt \cite{rootmap}.

Wie viel die Treffer wiegen, prüft eine Zufallsabschätzung nach einem vorab
festgelegten Verfahren. Der Katalog umfasst sechzehn Ziele, acht Formen und
ein Fenster von \(2\sigma\). Das Paar aus der Regel von M und der Regel von T
trifft dreimal, als einziges der sechsunddreißig möglichen Paare von
Zählregeln. Für ein zufälliges \(\gamma_L\) liegt die Wahrscheinlichkeit für
mindestens drei Treffer bei \(0{,}018\) bis \(0{,}023\). Über alle
sechsunddreißig Paare erwartet man daher \(0{,}6\) bis \(0{,}8\) Paare mit
mindestens drei Treffern. Nur wenn das Paar ohne Blick auf die Daten
festgelegt worden wäre, läge die Wahrscheinlichkeit bei etwa \(2\,\%\).
Solange das Paar gewählt und nicht hergeleitet ist, sind die
Treffer mit Zufall verträglich. Nachträglich angemerkt: Die beiden getroffenen
Ziele nahe \(0{,}30\) sind stark korreliert; zwei der drei Treffer sind daher
fast einer (Bemerkung v4.10.0-r30.5 \cite{chance,companion}).

Die Zahl selbst beruht darauf, zwei Frames verschiedene Zählregeln
zuzuordnen. Das ist gleichwertig mit einem frame-abhängigen Flächenquant.
Ein unitär wirkender Boost ändert das Spektrum des Flächenoperators aber
nicht, nur die Verteilung seiner Eigenwerte \cite{rovelli}. Ein
frame-abhängiges Quant bräuchte daher eine nicht-unitäre Frame-Abbildung.
Ein Vermittler, der alle fünf Gleichsetzungen aus einer Kopplung erzeugt,
ist im Companion nicht gefunden. Die Physik-Ausgabe von v5.3.0 behandelt dies in Abschnitt~7.7 \cite{physics}.
Auch ein Holst-Term in der Gravitationswirkung liefert die Abbildung nicht:
Lorentz-Transformationen, selbst ein komplexer Boost, lassen den
Barbero-Immirzi-Parameter fest, und nur eine Dualitätsdrehung der Krümmung um
einen Winkel \(\theta\) ändert ihn nach \(\arctan\gamma\mapsto\arctan\gamma+\theta\).
Kein im Companion ausgezeichneter Winkel bildet den Wert des Materie-Frames
auf den des zweiten Frames ab (Runde v4.10.0-r31 \cite{holst,layers}). Mehr noch: Jede Dualitätsdrehung,
reell oder komplex und in beliebiger Folge, erhält die Quadratsumme der
Einstein- und Holst-Koeffizienten. Für reelle Ergebnisse unterscheiden sich die
Flächenquanten beider Frames dann um
\(\sqrt{(1+\gamma_T^2)/(1+\gamma_M^2)}\approx0{,}98\), nicht um den bei gleicher
Gravitationskonstante nötigen Faktor \(\gamma_T/\gamma_M\).
Auch die nicht-unitäre Frame-Abbildung \(J_{\mathrm B}\) liefert ihn nicht (Runde v4.10.0-r52 \cite{layers}): Sie ist ein Automorphismus der Lorentz-Algebra, und eine
Konjugation, unitär oder nicht, lässt jedes Spektrum fest. Liest man die Fläche über die transformierte Simplizitätsbedingung ab, wird \(\gamma\)
quer zur Frame-Achse zu \(-1/\gamma\) und bleibt entlang der Achse; das ergibt weder \(\gamma_T\) noch \(\gamma_L\).
In den Einheiten \(r_j=e^{-2\pi\sqrt{j(j+1)}}\) ist jede Zählregel die Ähnlichkeitsgleichung einer selbstähnlichen Familie auf der Geraden mit
der Dimension \(\gamma<1\) (Runde v4.10.0-r53 \cite{layers}). Eine neuere, noch unbegutachtete Dimensionsformel über Entropieraten, die exakte Überlappungen zulässt \cite{openai148}, zeigt, was das
Gewicht 2 der Materie-Regel verlangt: Die beiden Extremzustände \(|j,\pm j\rangle\) müssen als verschiedene Abbildungen wirken. Fallen sie zusammen
oder wird nur einer zugelassen, sinkt der Wert auf \(\gamma_1=0{,}1616\), und die Regel des zweiten Frames fällt auf die Dimension null. Das Gewicht 2
verlangt also eine Auslese, die die beiden Orientierungen einer Punktierung unterscheidet, ohne eine auszuwählen; warum sie das tut, ist nicht
hergeleitet.

\subsection{Welche Messungen die Zahl prüfen}
\label{sec:data}

Aus der Zahl \(\beta\) und ihren Gleichsetzungen folgen Werte, die Messungen prüfen können. Der Companion hält für jede eine Vorab-Notiz
mit Prüfgröße, Prüfdatum und Entscheidungskriterium fest, bevor die entscheidenden Daten vorliegen \cite{predictions}. Die Tabelle
nennt die Werte und ihren Stand am 23.~November 2026. Ein Wert gilt erst als geprüft, wenn eine Messung nach dem festgelegten Stichtag
vorliegt; frühere Daten werden nur zum Vergleich angegeben.

\begin{center}
\small
\begin{tabular}{@{}p{.07\linewidth}p{.25\linewidth}p{.29\linewidth}p{.29\linewidth}@{}}
\hline
\textbf{Name} & \textbf{Vorhersage} & \textbf{entscheidende Messung} & \textbf{Vergleich heute}\\
\hline
H9 & \(\sin^2\theta_{12}=\beta^2/\cos^2\theta_{13}=0{,}3086\) & JUNO mit Fehler \(\le0{,}004\) & JUNO 2026: \(0{,}3036\pm0{,}0064\), Abstand \(+0{,}8\sigma\) \cite{juno2026}\\
H96 & \(\Omega_M=\beta^2=0{,}30185\) & erstes \(\Omega_m\) aus DESI DR3 (nur BAO) oder Euclid; widerlegt bei mehr als \(2\sigma\) & DESI DR2: \(0{,}2975\pm0{,}0086\), Abstand \(+0{,}5\sigma\) \cite{desi-dr2}\\
H102 & \(m_\tau=1777{,}134\)\,MeV aus der Restbiegung der Leptonphase & erste PDG-Ausgabe nach dem 13.~November 2026; widerlegt bei mehr als \(2\sigma\) & PDG 2026: \(1776{,}93\pm0{,}09\)\,MeV, Abstand \(+2{,}3\sigma\); zählt nicht, weil vor der Vorhersage \cite{pdg}\\
H112 & \(\sin^2\theta_{12}\cos^2\theta_{13}=\Omega_m\), ohne den Wert von \(\beta\) & die Messungen von H9 und H96 zusammen; widerlegt bei mindestens \(3\sigma\) & JUNO gegen DESI DR2: \(-0{,}1\sigma\)\\
H15 & \(m_W=80{,}3615\pm0{,}0018\)\,GeV & erste Weltmittelung nach dem 21.~Oktober 2026 & PDG 2026: \(80{,}3625\pm0{,}0077\)\,GeV; zählt nicht, weil vor der Vorhersage\\
\hline
\end{tabular}
\end{center}

Am 23.~November 2026 ist keiner dieser Tests entschieden. Drei Einschränkungen gehören zu jedem Treffer. Erstens hängen alle Werte bis
auf H112 an der Zuordnung der Zählregeln, also an OP1; die Zufallsabschätzung in Abschnitt~\ref{sec:one-number} gilt deshalb auch hier.
Zweitens sind H9 und H102 nach dem Sehen der älteren Daten gefunden worden; blind sind sie nur gegenüber künftigen Messungen. Drittens
prüft H112 eine schwächere Aussage: dass an beiden Stellen dieselbe Zahl steht, gleich welche. Er widerlegt erst ab einem Unterschied
von etwa \(0{,}022\) in \(\beta^2\). Weitere Vorhersagen zum Mischungswinkel \(\theta_{23}\) warten auf die nächste globale Auswertung
der Neutrinodaten \cite{predictions}.

\subsection{Welche Rechnung eine Annahme mitführt}

Für eine Nachrechnung wählt man eine Gleichung und schreibt ihre Eingaben auf. Beim kompakten gewöhnlichen Außenbeitrag sind das der lineare erhaltene mechanische Tensor, der Wellenfaktor des Profils und die retardierten Daten. Die gerichtete Auslese benutzt dieselbe Präparation, dazu den gewählten OP5-Antriebsoperator und die Kopplungsannahme in \eqref{eq:op5-match}. Die Gleichsetzung der Quellen ist der Übergang, durch den beide Wege Antworten auf dieselbe Materieänderung werden \cite{op5-motion,op5}.

QAM hält zu einer Formel ihre Darstellung, ihre Voraussetzungen und die noch offenen Übergänge für ihre Anwendung fest. Die im Companion protokollierten symbolischen Kontrollen gehören jeweils zu den durchgeführten Rechnungen. Eine Änderung der physikalischen Kopplung würde die Anwendung auf das gemeinsame Modell berühren. Die Erhaltungsrechnung für den mechanischen Tensor hat ihre eigenen, im Fallbeispiel sichtbaren Voraussetzungen \cite{qam,op5-motion}.

\paragraph{Wolfram-Rechenkontrolle.}
Für Fassung 0.1.0 liefen in Wolfram Language 15.0.1 insgesamt 80 exakte Prüfgruppen zum HC, zur Feldalgebra, zu Koide und zum geometrischen Vergleich. Die separate Dixon-Rechnung reproduzierte die Kommutantendimensionen sechs und null für den komplex-linearen und den antilinearen Anteil. Anhang~\ref{app:wolfram-output} stellt die Wolfram-Eingaben und ihre Originalausgaben einander gegenüber. Das Wolfram-Begleitarchiv enthält die Eingabeskripte, die protokollierten Ausgaben, den Prüfbericht und den geprüften Stand. Die Rechnungen kontrollieren die angegebenen mathematischen Identitäten und Testfälle unter ihren jeweiligen Voraussetzungen. Sie betreffen nur den Stand der Fassung 0.1.0. Teile, die danach hinzugekommen oder erweitert worden sind, decken sie nicht ab, darunter die Abschnitte~\ref{sec:octant-sign} und \ref{sec:one-number}, die Erweiterungen in Abschnitt~\ref{sec:open-points} und Anhang~\ref{sec:anhang-b}. Für die Herleitungen in Anhang~\ref{sec:anhang-b} gibt es eine eigene, enger gefasste Prüfung (Abschnitt~\ref{sec:anhang-b-pruefung}). Die neuen Zahlen in Abschnitt~\ref{sec:one-number} und das Ergebnis zur Auslese in Abschnitt~\ref{sec:open-points} belegt Anhang~\ref{app:new-calcs} mit den Originalausgaben; für die übrigen Aussagen gelten die Belege der jeweils genannten Companion-Runden \cite{companion,layers}.

Für die Vereinheitlichungsfrage ergibt sich so eine Folge von konkreten Aufgaben. Zuerst müssen die Zustände und ihre Ladungsdarstellung gewählt werden. Dann braucht die Massenauslese eine native Operatoridentifikation, einen Auswahlmechanismus und eine Skala. Schließlich muss sich dieselbe Materiekonstruktion in einer erhaltenen Quelle niederschlagen, deren Wirkung sich geometrisch beschreiben und prüfen lässt. Die Rechnungen geben an, an welchen Übergängen diese Schritte voneinander abhängen.

\subsection{Drehsinn und Turm der Hilbert-Shifts}
\label{sec:handedness-tower}

Nach Fassung 0.2.0 sind zwei Fragen in einer Reihe von Erkundungen und zwei Tests weiterverfolgt worden; Runde v4.10.0-r59 des Companion
fasst sie mit ihrem Status zusammen \cite{towerlayer}. Bis auf die beiden Tests H125 und H126 sind alle Schritte Erkundungen, deren Ergebnis
vor der Rechnung absehbar war; sie sind nicht blind und prüfen nichts.

\paragraph{Woher der Drehsinn kommt.}
Abschnitt~\ref{sec:octant-sign} und Runde v4.10.0-r58 führen das Vorzeichen des Oktantschritts auf das Produkt \(\tau\varepsilon\) aus der
Zeitrichtung der Auslese und dem Umlaufsinn der Bahn zurück. Fünf Erkundungen haben gefragt, ob ein weiterer Baustein diesen Drehsinn
festlegt: der Nullpunkt der Fläche einer Punktierung, die Restdrehung einer Scheibe aus drei Bögen als Holonomie, das Licht als Vermittler
zwischen den Frames, eine Störung des Lichtgewichts und die Drehrate einer Scheibe als Ursprung von \(\Lambda\). Jede Rechnung setzt den
Drehsinn voraus, keine legt ihn fest. Zwei Nebenergebnisse sind exakt: Mit dem halbklassischen Spektrum \(j+\tfrac12\) statt
\(\sqrt{j(j+1)}\) wird die Zählung geschlossen lösbar und gibt genau den Hopf-Wert \(\beta_{\rm sym}=1/\sqrt3\), verfehlt aber die Daten;
und der größte Formzyklus der Drei-Bogen-Scheibe dreht sie um genau \(-\pi\), sodass die Restdrehung \(\gamma_M\cdot\tfrac12\) wird, wenn man
\(\gamma_M\) als Kopplung einsetzt, nicht herleitet. Die Verknüpfung von \(\Lambda\) mit einer Drehrate ist nur eine Umbenennung, wie die
Epoche in Runde v4.10.0-r56. Am Ende ist die offene Wahl bei \(\varepsilon\) angekommen, also bei der Wahl der imaginären Einheit; über
\(\mathbb C\) ist das eine Konvention.

\paragraph{Ein relatives Vorzeichen (H125).}
Messbar ist dagegen ein Drehsinn zweier Sektoren relativ zueinander. Test H125 wurde vor der Rechnung festgelegt. Seine erste Stufe, eine
Strukturrechnung mit vorab erwartetem Ausgang, ist bestätigt: Die Familien-\(su(2)\) der Sedenionen (Runde v4.10.0-r14.4) ist auf den
Leptonen relativ zum Lorentz-Spin das komplex konjugierte Dublett, auf den Quarks nicht; die Bargmann-Invarianten sind
\(\tfrac14-\tfrac i4\) und \(\tfrac14+\tfrac i4\), in allen \(24+24\) zulässigen Zuordnungen. Die Wahl \(i\mapsto-i\) konjugiert beide Sektoren
zugleich und lässt das relative Vorzeichen stehen. Die zweite Stufe ist blind gegenüber künftigen Daten: Unter der Setzung, dass die
CP-Phase in beiden Sektoren auf dieselbe Weise aus der Familienstruktur entsteht, folgt \(\operatorname{sign}J_{\rm PMNS}=-\operatorname{sign}J_{\rm CKM}\),
also \(\sin\delta_{\rm PMNS}<0\). Entschieden wird mit der ersten Auswertung, die das Vorzeichen auf mindestens \(3\sigma\) festlegt,
spätestens Ende 2035. Ein einzelnes Vorzeichen hat die Grundwahrscheinlichkeit \(\tfrac12\); eine Bestätigung stützte vor allem die Setzung.
Runde v4.10.0-r65 schränkt diesen Befund ein: Das relative Vorzeichen hängt an einer weiteren benannten Setzung
(Abschnitt~\ref{sec:family-copies-koide}).

\paragraph{Der Turm der Hilbert-Shifts.}
Die zweite Frage war, ob der Feldwert eine Erzeugungsregel von einem Hilbert-Shift zum nächsten hat. Liest man den Feldwert auf der
Stufe \(n\) als Element der Cayley--Dickson-Algebra der Dimension \(2^n\), mit der Grundbahn \(\psi(t)=\exp(u\omega t)\psi_0\), dann ist die erste
Stufe eine reelle Schwingung, die zweite die Ellipse mit der Fläche als \(U(1)\)-Ladung, die dritte die Hopf-Sphäre mit der ersten
Holonomie und \(SU(2)\), die vierte die Oktonionen mit \(G_2\) und \(su(3)\) als Stabilisator einer Einheit. Danach kommt keine stetige
Symmetrie hinzu; die Kreisbahn schließt nicht mehr für jede Richtung, und Nullteiler treten auf. Neu ist bei den Sedenionen eine diskrete
Symmetrie \(S_3\), die drei Oktonionen-Kopien vertauscht, und die Zahl verträglicher Achsen folgt der Bott-Periode. Die Algebra ist bekannt;
neu ist nur die Zuordnung zum Turm, und sie ist eine Lesart. Rückwärts gelesen sind die meisten Wahlen oberhalb der ersten Stufe erzwungen
oder Konvention: Im Turm-Modell ist unter den erreichbaren Farbzahlen nur \(N_c=3\) anomaliefrei, ein Vorzeichenwechsel in der
Grundgleichung führt auf gespaltene statt kompakte Symmetrien, und eine CP-Phase gibt es erst ab drei Generationen. Mit einem Faktor
\(\exp(u\pi/4)\) in Kopienform erfüllt ein Zustand die Koide-Bedingung genau; die Phase \(\delta\) bleibt frei, und drei vollständige
Generationen passen in diese Algebra nicht. Ein Test, die Neutrino-Mischung aus der nicht-assoziativen Kopplung zweier Oktonion-Richtungen
zu gewinnen (H126), ist nach dem vorab festgelegten Kriterium nicht prüfbar, weil die Konstruktion überall entartet. Ein nichtlinearer
Zusatz in der Grundgleichung bringt eine Skala, ändert die Struktur aber nicht.

Keine dieser Rechnungen leitet \(\beta\), eine Masse, einen Mischungswinkel oder \(\Lambda\) her. Gesetzt bleiben der Rahmen der
Chiralität, die Unterscheidung von M und T, die Zahl der Generationen und jede Skala. OP1, OP3, OP4 und OP5 bleiben offen.

\subsection{Familiensymmetrie, Kopien und der Koide-Punkt}
\label{sec:family-copies-koide}

Nach Fassung 0.3.0 hat der Companion die Runden v4.10.0-r60 bis v4.10.0-r67 erhalten \cite{layers}. Die Runden r66 und r67 betreffen die
Zahl \(\beta\) und stehen in Abschnitt~\ref{sec:one-number}. Die Runden r60 bis r65 sind Sammelrunden und eine Korrekturrunde: Sie rechnen
nichts neu und halten Tests und Erkundungen mit dem Status fest, den die Ergebnis-Notizen nennen \cite{roundsr60r65}. Dieser Abschnitt gibt
sie verkürzt wieder; maßgeblich sind die Runden. Die erste Tabelle nennt die Tests mit vorab festgelegtem Kriterium. Wo ein Teil nicht blind
war, steht das in der dritten Spalte; die Runden führen es je Teilaussage.

\begin{center}
\small
\begin{tabular}{@{}p{.21\linewidth}p{.09\linewidth}p{.15\linewidth}p{.45\linewidth}@{}}
\hline
\textbf{Test} & \textbf{Runde} & \textbf{blind?} & \textbf{Ergebnis}\\
\hline
\raggedright H127: Brechung der Familiensymmetrie durch die Ladung & r60.2 & \raggedright vorab festgelegt & \raggedright Stufe A bestätigt: Die Brechung \(SO(3)\to A_4\) ist ein reiner Spin-3-Tensor. Stufe B nach dem Wortlaut widerlegt; sie war für jede zulässige Kopplung leer (r60.3).\tabularnewline
\raggedright H138--H147: Oktonion-Kopien in den Sedenionen, Rest-Generation & r62 & \raggedright teils nicht blind & \raggedright Widerlegt: H138a, H138b, H142b, H144c, nach dem Wortlaut H139a und H145c; die übrigen Teile bestätigt. Es gibt genau drei \(G_2\)-stabile Oktonion-Kopien, bedingt an die Definition der Kopie. Sie sind drei Lesarten einer Generation, nicht drei Generationen. Die Kopienstruktur liefert keine Massen.\tabularnewline
\raggedright H148--H156: der Koide-Punkt & r63 & \raggedright teils nicht blind & \raggedright Widerlegt: H150b in einem Teil, H155c, nach dem Wortlaut H150c und H155d; teilweise H151b und H156c; die übrigen Teile bestätigt, H153b bis auf eine Teilvorhersage. \(K=\tfrac23\) folgt aus einem Postulat mit gesetzter Skala, nicht aus der Algebra.\tabularnewline
\raggedright H128, H130: Familiensymmetrie und Träger in den Sedenionen & r64.1, r64.2 & \raggedright teils nicht blind & \raggedright H128a widerlegt (Dimension 22 statt der vorab genannten 17, numerischer Rang); H128b für die Richtung \(e_7\) bestätigt, für \(e_8\) widerlegt; H130a und H130b bestätigt.\tabularnewline
\raggedright H131, H134: Mischung aus den Kopien & r64.3 & \raggedright teils nicht blind & \raggedright H131a bestätigt, H131b widerlegt; H134a und H134c bestätigt, H134b widerlegt. Verträglichkeit mit den Daten, keine Vorhersage.\tabularnewline
\raggedright H129, H132, H133, H135, H136: rekursive Feldgleichung auf dem Turm & r64.4 & \raggedright teils nicht blind & \raggedright H129 und H133 nicht bestätigt; H132, H135 und H136 widerlegt. Der Ansatz legt keine Skala und keinen Wert fest.\tabularnewline
\raggedright H105: Flächenschritte und Echos & r64.5 & \raggedright vorab festgelegt, gegen künftige Daten & \raggedright offen; Frist 31.~Dezember 2028.\tabularnewline
\raggedright H162: Ladungsachse auf dem Sedenionen-Träger & r65.1 bis r65.3 & \raggedright a, b, d nicht blind & \raggedright a bis d bestätigt, c nur numerisch; Teil e ist eine Erkundung.\tabularnewline
\hline
\end{tabular}
\end{center}

\paragraph{Familiensymmetrie und Kopien.}
Die hermiteschen \(3\times3\)-Matrizen über den Oktonionen bilden die Albert-Algebra. Dort ist der Kommutant von \(G_2\) in
\(F_4\) eine Familien-\(SO(3)\); die Ladung bricht sie zu genau \(A_4\) (Runde v4.10.0-r60, Erkundungen k275 und k277, nicht blind). Die Gruppe \(A_4\),
die frühere Runden als Symmetrie der drei Plätze gesetzt hatten, ist dort also ein Rest. Bei exakter \(A_4\) sind die Quarkmassen entartet,
und die Quarkmischung ist die Einheitsmatrix; eine Mischung verlangt, dass \(A_4\) selbst bricht. Die Ausrichtung des Vakuums hängt an freien
Koeffizienten eines Potentials, sodass ein blinder Test derzeit nicht möglich ist. Die drei Plätze als Generationen zu lesen ist eine Lesart.
Runde v4.10.0-r62 fragt, ob die Sedenionen die Zahl der Generationen erklären. Die Abschlussbedingung einer Kopie ist kubisch und hat genau
drei Lösungen. Die drei Kopien haben dieselbe Farbe und tragen je eine Generation, aber keine gemeinsame Ladung, und zusammen nur 15 von 24
Zuständen. Auf der nächsten Stufe des Turms tritt bei vier Kopien in 9 von 126 Auswahlen eine vollständige Rest-Generation auf; ob ihre
Zustände physikalische Zustände sind, ist nicht gezeigt. Die Zahl der Generationen bleibt eine Eingabe.

\paragraph{Der Koide-Punkt.}
Liest man die Massen an den Amplituden ab (eine Setzung aus Runde v4.10.0-r59), dann gilt \(K=1/(3\cos^2\theta)\) mit einem Mischwinkel
\(\theta\), also \(K=\tfrac23\) genau bei \(\theta=\pi/4\) (Runde v4.10.0-r63). Dieser Winkel ist der Fixpunkt eines Tauschs \(J\), den der
Companion als Tausch von Materie- und Tachyon-Frame liest; jedes tauschsymmetrische Potential ist dort stationär. Hergeleitet sind diese
beiden Aussagen. Dass der Fixpunkt angenommen wird, folgt nur aus dem Postulat, dass die beiden Frames gleichberechtigt sind, und mit
gesetzter Skala; der Autor hat das Postulat am 26.~November 2026 als Arbeitsannahme angenommen, hergeleitet ist es nicht. Der Tausch \(J\)
ist antilinear und weder die Tachyon-Abbildung \(z\mapsto iz\) noch die Zeitumkehr \(z\mapsto\bar z\) (Bemerkung v4.10.0-r64.7). Die Phase
\(\delta=2/9\) bleibt angepasst. Ohne gesetzte Skala verschiebt der Fluss den Wert abhängig vom Start. Der gemessene Wert
\(K=0{,}6666645\pm0{,}0000051\) liegt \(0{,}43\sigma\) unter \(\tfrac23\). Runde r63 führt daraus eine bedingte, offene Vorhersage
\(m_\tau=1776{,}969\)\,MeV, falls \(K=\tfrac23\) für die Polmassen exakt gilt; sie ist nicht die Vorhersage H102 aus
Abschnitt~\ref{sec:data}.

\paragraph{Ladungsachse und das relative Vorzeichen.}
Abschnitt~\ref{sec:handedness-tower} nennt das relative CP-Vorzeichen zwischen Leptonen und Quarks (H125). Es war an eine Ladungsachse
gebunden, die ein Platzhalter des Companion ist (CT-003); Test H130 stützt diese Achse nicht (Runde v4.10.0-r64). Test H162 zeigt, dass das
Spektrum der Ladung, die Konstruktion mit zwei Generationen, die Familien-\(su(2)\) und die erste Stufe von H125 für jede Achse gelten
(Runde v4.10.0-r65). Eine Erkundung derselben Runde (k361, nicht blind, nur numerisch) findet zwei Typen der Ladung, mit Links- oder mit
Rechtsmultiplikation, und schränkt den Befund aus Abschnitt~\ref{sec:handedness-tower} ein: Das relative Vorzeichen hängt nicht an der Wahl
\(i\mapsto-i\), nicht an der Achse und nicht am Ladungstyp, wohl aber an der Setzung, dass dieselbe Seite der Multiplikation die Familie in
beiden Sektoren trägt. Sie liegt nahe, folgt aber nicht aus Ladung und Farbe allein. Die Vorhersage \(\sin\delta_{\rm PMNS}<0\) bleibt offen
und hängt jetzt an zwei benannten Setzungen. Der Träger enthält zwei Generationen, nicht drei.

\paragraph{Korrekturen und Nachrechnung.}
Runde v4.10.0-r61 berichtigt zwei Formeln früherer Runden (Erkundung k320, nicht blind): die Normierung der Vakuumdichte,
\(\rho_\Lambda=\Lambda/\kappa\), und die Wurzel in \(\beta=\sqrt{(\gamma_L-1)/(\gamma_L+1)}\). Keine Zahl, kein Testergebnis und kein Status
ändert sich; dieser Artikel hat keine der beiden Formeln zitiert. Runde r64 präzisiert außerdem den Geltungsbereich einiger älterer Aussagen,
ebenfalls ohne Statusänderung. Nach der Annahme des Zählpostulats hat eine Erkundung (k358, kein Test; Runde v4.10.0-r68) die Zahlen, die von \(\beta\)
abhängen, mit dem exakten Wert nachgerechnet \cite{recalc}: Kein Urteil ändert sich; am schärfsten reagiert \(m_W\), dessen Abstand zu
PDG~2026 von \(-0{,}44\sigma\) beim zuvor geführten, aus gerundeten Eingaben gewonnenen Wert auf \(-0{,}13\sigma\) geht, und \(m_\tau\) aus
H102 bleibt bei \(+2{,}3\sigma\).

Keine dieser Runden leitet \(\beta\), eine Masse, einen Mischungswinkel oder \(\Lambda\) her. Gesetzt bleiben die Zahl der Generationen,
die Definition der Kopie, die Ladungsachse, die Seite der Multiplikation, die Auslese an den Amplituden und jede Skala; angepasst bleibt
\(\delta\). OP1 bis OP5, C5 und C6/T2 bleiben offen, und kein Guard ist aufgehoben.

\subsection{Was nach dieser Fassung offen bleibt}
\label{sec:open-points}

Der Companion führt seine offenen Probleme unter festen Namen. Ein Guard ist dort eine ausdrücklich gesperrte Aussage: Sie darf nicht
als bewiesen verwendet werden, solange ihr Nachweis fehlt. Die folgende Tabelle nennt die Probleme mit einer Zeile; keines ist
geschlossen, und kein Guard ist aufgehoben \cite{companion,layers}.

\begin{center}
\small
\begin{tabular}{@{}p{.12\linewidth}p{.55\linewidth}p{.27\linewidth}@{}}
\hline
\textbf{Name} & \textbf{Frage} & \textbf{im Artikel}\\
\hline
OP1 & Woher kommt die Zahl \(\beta=0{,}5494\), also warum zählen die Frames M und T die Horizontzustände nach verschiedenen Regeln? & Abschnitt~\ref{sec:one-number}\\
OP2 & Existiert der Dirac-artige Operator auf dem antisymmetrischen Sektor als selbstadjungierter Operator, ohne die zurückgezogene Riemann-Arbeit, und liegt sein Spektralversatz bei \(\beta\)? & nur als Gleichsetzung in Abschnitt~\ref{sec:one-number}\\
OP3 & Welche Übergänge der Hülle sind zulässig, woher kommen drei Generationen, und was legt Phase und Vorzeichen der Leptonphase fest? & Abschnitte~\ref{sec:charged-multiplicity}, \ref{sec:octant-sign}, \ref{sec:handedness-tower} und \ref{sec:family-copies-koide}\\
OP4 & Woher kommen die Mischungswinkel der Neutrinos und der Quarks? & Gleichsetzung in Abschnitt~\ref{sec:one-number}; relatives CP-Vorzeichen in den Abschnitten~\ref{sec:handedness-tower} und \ref{sec:family-copies-koide}\\
OP5 & Wie fügen sich Vakuumbilanz und kosmologische Konstante in den Rahmen, und welche Materiequelle treibt welches geometrische Feld? & Fallstudie in Abschnitt~\ref{sec:op5-coupling}, Anhang~\ref{sec:anhang-b-annahmen}\\
C5 & Kann die verbleibende Obstruktion in endlich vielen Schritten exakt null werden, statt nur unter eine Detektorschwelle zu fallen? & nicht behandelt\\
C6/T2 & Schließen alle Auslesungen zugleich, global und nicht nur lokal? & nicht behandelt; durch Guards gesperrt\\
\hline
\end{tabular}
\end{center}

Drei Fragen tragen die weitere Arbeit. Sie sind hier mit ihrem Stand zusammengefasst; die Herleitungen dazu stehen im Companion.

Wie sie zusammenhängen, ordnet Runde v4.10.0-r55 des Companion \cite{crossmap,layers}. Alle bedingten Zahlenwerte hängen an der einen Zahl aus
Abschnitt~\ref{sec:one-number}. Quer zu den einzelnen Problemen liegen zwei Lücken. Die erste ist die Auslese: Mehrere Stellen brauchen ein
Bezugssystem, eine Achse und eine Zeitorientierung, die der Companion nicht festlegt. Sie zerfällt in die Frage nach der Orientierung, die
beim Vorzeichen des Oktantschritts und bei der Zählung der Horizontzustände auftritt, und die Frage nach Bezugssystem und Achse, die die
Beobachterlesart des Materie-Frames und die gerichtete OP5-Auslese betrifft. Die zweite Lücke ist eine gemeinsame Dynamik mit eigener Skala,
die bisher fehlt. Daneben hat jedes Problem einen eigenen Rest, der von der Zahl nicht abhängt. Die erste der drei Fragen unten gehört zur
Orientierung, die zweite zu einem solchen Rest, die dritte zu Bezugssystem und Achse und zur fehlenden Dynamik.

Für einen dieser Reste liegt inzwischen ein Negativ-Schluss vor. Der eigene Rest von OP5 ist der Absolutwert der kosmologischen
Konstante. Runde v4.10.0-r56 zeigt für eine ausdrücklich genannte Klasse von Konstruktionen des Companion (Frame-Abbildungen und
Auslesungen, die Wirkungen W1/W2, den semiklassischen Weg, das OP5-Vergleichsmodell und die Kaskade), dass eine additive Konstante aus
jeder Bewegungsgleichung herausfällt und mit spurfreien Gleichungen eine Integrationskonstante ist, dass diese Klasse keine eigene Länge
enthält und dass die Auslesung \(\Omega_M=\beta^2\) für jeden Wert von \(\Lambda\) zu genau einer Epoche gilt; sie gibt nur
\(\Lambda=3(1-\beta^2)H^2\) an dieser Epoche \cite{lambdainput}. Der Wert von \(\Lambda\) ist dort also eine Eingabe, gleichwertig die Wahl
der Epoche \glqq heute\grqq{}. Das ist ein Ausschluss für den vorhandenen Bestand, nicht für jede Erweiterung; OP5 bleibt offen, und die
Vorhersage \(\Omega_M=\beta^2\) wartet auf die Messung (Abschnitt~\ref{sec:one-number}).
Der andere Weg wäre eine Dynamik mit \(\Omega_M=\beta^2\) als Attraktor. Die einzige Dynamik des Companion für das Verhältnis, die Kaskade,
in der Vakuum mit der Rate \(\beta^2\varepsilon H\) in Materie zerfällt, hat den Attraktor \(\varepsilon\beta^2/3\); gleich \(\beta^2\) ist er genau für
\(\varepsilon=3\). Gegen die BAO-Daten von DESI DR2 ist dieses Modell nach einem vorab festgelegten, nicht blinden Vergleich nicht
ausgeschlossen, verlangt aber heute \(\Omega_M\approx0{,}50\) und passt deutlich schlechter als \(\Lambda\)CDM (Runde v4.10.0-r57 \cite{layers}).
Der Wert \(\beta^2\) läge dann in der Zukunft; warum er heute gemessen wird, erklärt der Attraktor nicht.

\paragraph{Ein nativer physikalischer Operator.}
Die Massenauslese braucht einen Operator, der aus dem gemeinsamen Träger folgt und nicht nachträglich gewählt ist. Bisher ist er nicht
gefunden. Abschnitt~\ref{sec:octant-sign} zeigt, dass der vorhandene Kompensator auf den geschlossenen ungeraden Zuständen verschwindet
und dass kein Baustein des Companion das Vorzeichen des Oktantschritts festlegt; es ist eine Zeitorientierung der Auslese. Dieselbe
Orientierung tritt bei der Zählung der Horizontzustände wieder auf: Das Gewicht 2 der Materie-Regel verlangt eine Auslese, die die beiden
Orientierungen einer Punktierung unterscheidet, ohne eine auszuwählen (Runde v4.10.0-r53 \cite{layers}). Ein nativer Operator müsste beides zugleich
erklären, das Vorzeichen im Einzelfall und die Unterscheidung in der Zählung. Eine Auslese in drei zeitlich geordneten Fenstern, die
Bargmann-Invariante des Oktant-Tripels auf einer Punktierung vom Spin \(j\), leistet das nicht (Runde v4.10.0-r58 \cite{layers}, Anhang~\ref{app:new-temporal}):
Sie gibt beiden Orientierungen denselben Betrag und die Werte \(2^{-3j}e^{\pm ij\pi/2}\), unterscheidet sie also nur für \(j\), die keine geraden
ganzen Zahlen sind. Als Zählregel gelesen, mit Gewicht 1 auf den geraden \(j\), ergibt sie \(\gamma_L'=1{,}838\) statt \(1{,}865\) und verfehlt
\(m_W\) um rund \(20\sigma\). Vorzeichen und Zählung kommen damit nicht aus derselben Auslese dieser Form. Abschnitt
\ref{sec:one-number} zeigt, dass alle bedingten Zahlenwerte an einer Zahl hängen, die über gesetzte Gleichsetzungen eingeht.

\paragraph{Eine hergeleitete Zuordnung der Eichladungen.}
Die Ladungszuordnung im Dixon-Vergleich ist konstruiert, nicht abgeleitet. Abschnitt~\ref{sec:charged-multiplicity} grenzt ein, wo ein
Generationsindex sitzen kann: nicht in der Farbe, sondern in einer farbneutralen Multiplizität, deren Zahl \(n\) eine Eingabe bleibt.

\paragraph{Ein unabhängig begründetes geometrisches Quellgesetz.}
Für den gewöhnlichen Teil ist es vorhanden: Der präparierte Beitrag in Abschnitt~\ref{sec:op5-coupling} folgt allein aus der
linearisierten Einstein-Gleichung und ist außerhalb seines Trägers kompakt (Anhang~\ref{sec:anhang-b-kompakt}). Der gerichtete Kontrast
braucht zusätzlich die Gleichsetzung der Higgsquelle mit dem Antrieb des Vergleichsfelds; am Higgs-Sektor der Wirkung W2 folgt sie nur für
den skalaren Kanal, und dann fehlt der gerichtete Vermittler, an dem der Kontrast hängt (Runde v4.10.0-r33 \cite{layers}). Die Suche nach
diesem Vermittler in den Runden v4.10.0-r34 bis v4.10.0-r54 trägt nur teilweise: Gesetzt bleiben ein massiver Spin-2-Vermittler, seine
Kopplung an den Higgsradius und eine aus Gründen der Kausalität verlangsamte Ausbreitung, und Bezugssystem und Achse sind bei OP5 eine
Setzung des Detektors, keine Wahl des Feldes \cite{layers}. Mit dem lichtschnellen Higgsradius von W2 ist die gerichtete Kopplung für keine
Wahl der Vermittler kausal verträglich, und im Vakuum begrenzen Laborgrenzen der Lorentz-Verletzung sie auf \(g_m^2\lesssim10^{-26}\); ein
beobachtbarer Kontrast verlangt ein Material. Offen ist damit, ob es einen solchen Vermittler als eigenes Feld gibt; die einzelnen Schritte
stehen in Anhang~\ref{sec:anhang-b-vermittler}.

Keiner dieser Punkte ist durch die vorliegende Fassung geschlossen. Sie legt offen, an welcher Stelle eine Herleitung ansetzen müsste.

